\phantomsection
\chapter*{Aku Hanya Ingin di Anggap Ada}
\addcontentsline{toc}{section}{Aku Hanya Ingin di Anggap Ada --- Ricky Juliana Sodikin}
\penulis{Ricky Juliana Sodikin}
\awalcerpen{D}{i desa Welasari} hiduplah satu keluarga yang cukup bahagia dan rukun, didalamnya terdapat sepasang kekasih bernama Dirgantara dan Reyna, mereka di buahi 2 buah anak, anak pertama mereka bernama Samudra, sedangkan anak ke 2 mereka bernama Langit. Kehidupan mereka sangatlah baik, hingga suatu ketika.


“Langit, kesini sebentar nak”  Ucap Dirgantara memanggil anaknya.


“Iya ayah?” Langit menghampiri Ayahnya dan menjawab panggilan itu.


“Ayah punya hadiah buat kamu loh, nis laptop yang kamu minta ke ayah kemarin, di jaga ya, jangan di rusakkin” Ucap ayahnya sembari tersenyum


Ketika Dirgantara dan anaknya Langit mengobrol di ruang keluarga, Samudra ternyata melihat Ayahnya memberikan hadiah kepada adiknya, dirinya pun menghampiri Ayahnya dengan maksud untuk menanyakan apakah ada hadiah untuk dirinya.


“Ayah, hadiah buat Samudra mamna?” Ucap Samudra dengan rasa excited yang tinggi


“Minta hadiah? Belajar, bukannya minta minta, sana udah ke kamar” Ucap Dirgantara dengan nada tinggi.


Ibunya, Reyna menghampiri mereka di ruang tamu.


“Udah bang sana ke kamar, jangan manja udah gede kamu itu” Ucap Reyna.


Samudra pun pergi ke kamar dengan rasa yang sedih.


Keesokan harinya keluarga tersebut berkumpul di sebuah meja makan, ibunya membuat sebuah bolu untuk Langit sebagai hadiah darinya karena Langit mendapatkan ranking 3 di kelasnya.


“Wiih bolu! Makasih Ibu!” Ucap Langit.


“Iya sayang, sama sama” Ucap Reyna.


“Waah, Pasti enak ini bolu buatan Ibu kamu, di habisin ya Dek” Ucap Dirgantara.


Samudra yang melihat Langit di rayakan dengan di berikan sebuah bolu merasa iri, dengan rasa keberanian dirinya berbicara ke Ibunya.


“A abang ada tidak hadiah dari ibu?” Ucap Samudra dengan gugup.


“Ini khusus buat Langit, kamu kapan kapan aja ya?. udah yuk makan ibu udah masak cape cape ini” Ucap Reyna.


Samudra merasa sedih, karena dia merasa bahwa dirinya hanya menjadi bayangan di antara keluarganya, dia merasa bahwa semua perhatian hanya untuk Langit. Samudra pun pergi dari hadapan mereka.


Samudra pergi ke kamarnya dan air matapun tak dapat dia tahan.


“Kenapa aku yang di perhatiin pake cara yang beda? Kenapa selalu Langit yang dapet perhatian penuh dari kalian? Ga ada rasa sayang ke Samudra kah?” Samudra menangis se menjadi menjadinya.


Setelah beberapa kejadian yang tidak berpihak ke Samudra, Samudra pun menjadi seseorang yang pendiam, dirinya menjadi pribadi yang tertutup kepada orang tuanya, ketika dirinya bertemu dengan keluarganya, tak ada sepatah katapun yang keluar dari mulutnya untuk keluarganya.


Keesokan harinya, disaat semua orang akan beraktivitas, Samudra menghampiri kedua Orang tuanya di ruang tamu dengan badan yang lumayan panas.


“A ayah, I Ibu… Samudra cuma pengen nanya, sebenernya samudra ini di anggap ada ga sih di keluarga ini? Kayaknya cuma Samudra yang ga dapet kasih sayang kalian disini, Samudra bingung.. apa samudra emang ada salah yang bikin samudra seperti ini? Samudra minta maaf untuk semuanya.. Samudra rasa Samudra disini hanya seperti bayangan, Samudra cuma pengen di anggap ada..” Ucap Samudra sembari menangis.


Kedua orang tuanya yang mendengar perkataan Samudra pun seakan tersadar, bahwa mereka selama ini terlalu menyayangi salah satu anaknya, ruang tamu hening dengan di iringi air dari keran menetes dari toilet.


“I Ibu minta maaf ya Samudra.. Ibu sadar peran Ibu tidak ada untuk kamu selama ini, Ibu terlalu memberikan kasih sayang, perhatian ke Langit, Ibu minta maaf…” Ucap Ibunya sembari menunduk.


“Ayah juga minta maaf ya nak, atas semua kekurangan Ayah ke kamu.. Ayah benar benar minta maaf Nak..” Ucap Ayahnya.


Samudra yang mendengar itu merasa senang, karena akhirnya dari sekian lama dirinya di anggap ada kembali oleh orang tuanya. namun setelah mendengar itu, Samudra terjatuh dan pingsan akibat Badan yang mengalami demam tinggi. dirinya pun di bawa ke rumah sakit oleh kedua orang tuanya.


Sejak hari itu, keluarga mereka berubah, Samudra yang dulunya selalu mendapatkan kasih sayang berbeda dari orang tuanya, kini dirinya mendapatkan kasih sayang penuh dari kedua orang tuanya. Keluarga mereka kini hidup dengan damai, canda tawa mereka buat, momen indah tercipta di antara mereka. 

\jedahias
